\documentclass[10pt,a4paper]{amsart}
\usepackage[margin=19mm]{geometry}
\usepackage{amsmath,amssymb,mathtools}
\usepackage[scr=rsfs,cal=euler]{mathalfa}
\usepackage{xfrac}
\usepackage[unicode,hidelinks,bookmarks=false]{hyperref}
\newcommand{\ie}{i.e.\ }
\newcommand{\cf}{cf.\ }
\newcommand{\resp}{respectively\ }
\newcommand{\Subsection}{\subsection}
\providecommand{\nobreakdash}{\nobreak\hbox{-}}
\setlength{\emergencystretch}{2em}
\multlinegap0pt
\usepackage{bm}
\usepackage{bbm}
\usepackage{dsfont}
\usepackage{xspace}
\usepackage{cancel}
\usepackage{mathtools}
\usepackage{amsthm}
\makeatletter
\@ifundefined{theorem}{\newtheorem{theorem}{Theorem}}{}
\makeatother
\title[Pr\"ufer Transformation and Spectral Analysis for a Sturm--L]{Pr\"ufer Transformation and Spectral Analysis for a Sturm--Liouville-Type Equation}
\author{Shalmali Bandyopadhyay, F. Ay\c{c}a \c{C}etinkaya, Tom Cuchta}
\date{Source: arXiv:2508.12369v2 (CC BY 4.0)}
\begin{document}
\maketitle

\noindent\emph{Source and licence.} Pr\"ufer Transformation and Spectral Analysis for a Sturm--Liouville-Type Equation. Version arXiv:2508.12369v2 (CC BY 4.0), distributed under the Creative Commons Attribution 4.0 International licence, as recorded in the arXiv licence metadata for this exact version and shown on the abstract page https://arxiv.org/abs/2508.12369v2. Full rights notice: licenses/arxiv-2508.12369-CC-BY-4.0.txt.\par\smallskip

\noindent\textit{Editorial scope.} Counted unit: the Theorem whose source label is \texttt{th:upper-bound} in the article (source0.tex lines 628--642, source label \texttt{th:upper-bound}), delivered together with the article's own complete original proof of it (source0.tex lines 644--769, one \texttt{proof} environment of 126 lines). No mathematics is written, repaired or re-derived: statement and proof are the article's own text line for line. Required context retained uncounted: 14 (lines 434--436); RayleightQuotient (lines 492--494). Every definition, display and label of those lines is retained unchanged. Omitted from this package and not redistributed: 1--433, 437--491, 495--627, 643--643, 770--990. Nothing inside a retained excerpt is omitted or altered: every retained body is delivered exactly as the article's lines are written, so no omission receipt of any kind accompanies this package. Citations: the retained excerpts carry no citation command at all, so no bibliography entry of the article is transcribed and no reference is redistributed. Reference adaptation: none was needed and none was made; every reference inside the delivered unit resolves inside it, so no unresolved reference or citation is produced. Printed slips kept literal: no printed word, symbol, hypothesis or number of the counted statement or of its proof was altered. Layout and class: the article's own class is \texttt{amsart} with packages including inputenc, fontenc, lmodern, microtype, float, placeins, geometry, amsmath, amssymb, amsfonts, mathtools, mathrsfs, bm, amsthm; the wrapper is the lane's pinned \texttt{amsart} editorial wrapper at 10pt on A4, which restates the theorem-like environments the retained text needs and adds the article's own macro definitions that the retained text needs (none), with extra wrapper packages (bm, bbm, dsfont, xspace, cancel, mathtools). The delivered numbering is the unit's own. Assets: no retained span contains a figure environment, a graphics-inclusion command, a PDF-page inclusion, a table or a TikZ picture; no image file is copied into this package, the package declares no asset dependency and no image file is redistributed with it. Licence: version arXiv:2508.12369v2 of this article is distributed under the Creative Commons Attribution 4.0 International licence, as recorded in the arXiv licence metadata for this exact version and shown on the abstract page https://arxiv.org/abs/2508.12369v2. A full rights notice is delivered at licenses/arxiv-2508.12369-CC-BY-4.0.txt. Characters: every retained excerpt is pure ASCII, so the delivered engine, XeLaTeX with its default Latin Modern text font, reports no missing character.
\begin{equation} \label{14}
-\left(p(x)\left(y'(x) + s(x)y(x)\right)\right)' + s(x)p(x)\left(y'(x) + s(x)y(x)\right) + q(x)y(x) = \lambda \omega(x)y(x).
\end{equation}

\begin{equation}\label{RayleightQuotient}
\mathcal{R}[v] = \frac{\displaystyle\int_a^b \tilde{p}(x)(v'(x))^2\,\mathrm{d}x + \int_a^b Q(x)v^2(x)\,\mathrm{d}x}{\displaystyle\int_a^b \tilde{\omega}(x)v^2(x)\,\mathrm{d}x},
\end{equation}

\begin{theorem}[Upper bound on the $n^{\text{th}}$ eigenvalue]
\label{th:upper-bound}
Let $\tilde{p}(x) = p(x)e^{-2\bar{s}(x)}$ and $\tilde{\omega}(x) = \omega(x)e^{-2\bar{s}(x)}$, where $\bar{s}(x) = \int_a^x s(t)\,\mathrm{d}t$. Suppose $h \in C([a,b])$ satisfies $h(x) \geq \tilde{p}(x)$ for all $x \in [a,b]$, and define
\[
D := \min_{x \in [a,b]} \tilde{\omega}(x)h(x) > 0.
\]
Let
\[
M := \max_{x \in [a,b]} \frac{q(x)}{\omega(x)}.
\]
Then the $n^{\text{th}}$ eigenvalue of~\eqref{14} satisfies
\[
\lambda_n \leq \frac{n^2\pi^2}{D\left(\displaystyle\int_a^b \frac{1}{h(x)}\,\mathrm{d}x\right)^2} + M, \qquad n = 1, 2, 3, \ldots
\]
\end{theorem}

\begin{proof}
We again write the Rayleigh quotient~\eqref{RayleightQuotient} in the form
\[
\mathcal{R}[v]
= \frac{\displaystyle\int_a^b \tilde{p}\,(v')^2\,\mathrm{d}x}
       {\displaystyle\int_a^b \tilde{\omega}\,v^2\,\mathrm{d}x}
+ \frac{\displaystyle\int_a^b Q\,v^2\,\mathrm{d}x}
       {\displaystyle\int_a^b \tilde{\omega}\,v^2\,\mathrm{d}x}.
\]
Since
\[
M=\max_{x\in[a,b]}\frac{Q(x)}{\tilde{\omega}(x)}
   =\max_{x\in[a,b]}\frac{q(x)}{\omega(x)},
\]
we have $Q(x)\le M\,\tilde{\omega}(x)$ for all $x\in[a,b]$. Consequently,
\[
\frac{\displaystyle\int_a^b Q\,v^2\,\mathrm{d}x}
     {\displaystyle\int_a^b \tilde{\omega}\,v^2\,\mathrm{d}x}
\le M.
\]
Moreover, since $h\ge\tilde{p}$ on $[a,b]$, we have
\[
\frac{\displaystyle\int_a^b \tilde{p}\,(v')^2\,\mathrm{d}x}
     {\displaystyle\int_a^b \tilde{\omega}\,v^2\,\mathrm{d}x}
\le
\frac{\displaystyle\int_a^b h\,(v')^2\,\mathrm{d}x}
     {\displaystyle\int_a^b \tilde{\omega}\,v^2\,\mathrm{d}x}.
\]

We now consider the self-adjoint comparison problem
\begin{equation}\label{comparison_upper}
-(h(x)\psi'(x))' = \mu\,\frac{\psi(x)}{h(x)},
\qquad
\psi(a)=\psi(b)=0.
\end{equation}
Multiplying by $\psi$ and integrating by parts yields the Rayleigh quotient
\[
\tilde{\mathcal{R}}[\psi]
=
\frac{\displaystyle\int_a^b h\,(\psi')^2\,\mathrm{d}x}
     {\displaystyle\int_a^b \frac{\psi^2}{h}\,\mathrm{d}x}.
\]
Applying the Liouville substitution
\[
u=\int_a^x \frac{\mathrm{d}t}{h(t)}
\]
reduces this problem to the Dirichlet eigenvalue problem for the second derivative on the interval
\(
(0,L)
\),
where
\(
L=\int_a^b \frac{\mathrm{d}t}{h(t)}.
\)
The corresponding eigenvalues are
\[
\mu_n
=
\frac{n^2\pi^2}{L^2}
=
\frac{n^2\pi^2}
{\left(\displaystyle\int_a^b \frac{1}{h(x)}\,\mathrm{d}x\right)^2},
\qquad n=1,2,3,\ldots
\]
The hypothesis $\tilde{\omega}h\ge D$ implies
\[
\tilde{\omega}\ge \frac{D}{h} \quad \text{for all } x\in[a,b].
\]
Therefore,
\[
\int_a^b \tilde{\omega}\,v^2\,\mathrm{d}x
\ge
D\int_a^b \frac{v^2}{h}\,\mathrm{d}x,
\]
and it follows that
\[
\frac{\displaystyle\int_a^b h\,(v')^2\,\mathrm{d}x}
     {\displaystyle\int_a^b \tilde{\omega}\,v^2\,\mathrm{d}x}
\le
\frac{1}{D}
\frac{\displaystyle\int_a^b h\,(v')^2\,\mathrm{d}x}
     {\displaystyle\int_a^b \frac{v^2}{h}\,\mathrm{d}x}=\dfrac{\tilde{\mathcal{R}}[v]}{D}
\]
where $\tilde{\mathcal{R}}[v]$ is the Rayleigh quotient for the comparison problem~\eqref{comparison_upper}. 

Let $\psi_1,\dots,\psi_n$ denote the first $n$ eigenfunctions of ~\eqref{comparison_upper} and define
\[
V=\mathrm{span}\{\psi_1,\dots,\psi_n\}.
\]
Then $\dim V=n$, and for every $v\in V\setminus\{0\}$ we have
\(
\tilde{\mathcal{R}}[v]\le\mu_n
\).
Combining the previous inequalities, it follows that
\[
\mathcal{R}[v]
\le
\frac{1}{D}\,\tilde{\mathcal{R}}[v] + M
\le
\frac{\mu_n}{D} + M,
\qquad v\in V\setminus\{0\}.
\]
Taking the maximum over $v\in V\setminus\{0\}$ and invoking the minimax characterization of eigenvalues yields
\[
\lambda_n
=
\min_{\substack{S\subset H_0^1(a,b)\\ \dim S=n}}
\max_{\substack{v\in S\\ v\neq 0}}
\mathcal{R}[v]
\le
\max_{\substack{v\in V\\ v\neq 0}}
\mathcal{R}[v]
\le
\frac{\mu_n}{D} + M.
\]
Substituting the expression for $\mu_n$ gives
\[
\lambda_n
\le
\frac{n^2\pi^2}
     {D\left(\displaystyle\int_a^b \frac{1}{h(x)}\,\mathrm{d}x\right)^2}
+ M,
\]
which completes the proof.
\qedhere
\end{proof}

\end{document}
